\documentclass[main.tex]{subfiles}


\begin{document}

%from pagenumber to up
% \phantomsection 
% \hypertarget{toc}{}

 

 

 
   

\section{Магнитное поле катушки с током}


Пусть имеется цилиндрическая \emph{катушка}, состоящая из нескольких витков провода, намотанного по винтовой линии (\emph{соленоид}). На рис.\,\ref{fig:p151} показана картина линий магнитного поля катушки с током.


\pgfkeys{%
/pgf/decoration/.cd,
3d coil color/.store in=\TDCoilColor, 
3d coil color/.initial=black,
3d coil color=black,
3d coil width/.store in=\TDCoilWidth, 
3d coil width/.initial=0.4pt,
3d coil width=0.4pt,
3d coil dist/.store in=\TDCoilDist, 
3d coil dist/.initial=0.6pt,
3d coil dist=0.6pt,
3d coil opacity/.store in=\TDCoilOpacity, 
3d coil opacity/.initial=1,
3d coil opacity=1
}

\makeatletter % https://tex.stackexchange.com/a/219088/121799
\tikzset{get stroke color/.code={%
    \expandafter\global% Jump over, now we have \global
    \expandafter\let% Jump over now we have \global\let
    \expandafter\pgfsavedstrokecolor% Jump we have \global\let\pgf...
    \csname\string\color@pgfstrokecolor\endcsname% Finally expand this and put it at the end 
    },                                           % \global\let\pgf...{} in expanded form 
    restore stroke color/.code={\pgf@setstrokecolor#1},
}
\def\pgfpoint@onthreedcoil#1#2#3{%
  \pgf@x=#1\pgfdecorationsegmentamplitude%
  \pgf@x=\pgfdecorationsegmentaspect\pgf@x%
  \pgf@y=#2\pgfdecorationsegmentamplitude%
  \pgf@xa=0.083333333333\pgfdecorationsegmentlength%
  \advance\pgf@x by#3\pgf@xa%
  \advance\pgf@x by-\initialoffset pt%
}

% coil decoration
%
% Parameters: \pgfdecorationsegmentamplitude, \pgfdecorationsegmentlength,

\pgfdeclaredecoration{3d coil}{initial}
{ 
    \state{initial}[width=0.25*\pgfdecorationsegmentlength+\pgfdecorationsegmentaspect*\pgfdecorationsegmentamplitude,
    next state=coil, persistent precomputation={
    \tikzset{get stroke color}
    \pgfmathsetmacro{\initialoffset}{0.25*\pgfdecorationsegmentlength+\pgfdecorationsegmentaspect*\pgfdecorationsegmentamplitude}
  }]
  {%
    \pgfpathmoveto{\pgfpointorigin} 
    \pgfsetstrokecolor{\TDCoilColor}
    \pgfsetstrokeopacity{\TDCoilOpacity}
    \pgfsetlinewidth{1.5*\TDCoilWidth} 
    \pgfpathcurveto
    {\pgfpoint@oncoil{0    }{ 0.555}{1}}
    {\pgfpoint@oncoil{0.445}{ 1    }{2}}
    {\pgfpoint@oncoil{1    }{ 1    }{3}}
    \pgfusepath{stroke}
    \pgfcoordinate{TD@coilast}{\pgfpoint@oncoil{1    }{ 1    }{3}} 
    }
  \state{coil}[switch if less than=%
    1.25\pgfdecorationsegmentlength+%
    \pgfdecorationsegmentaspect\pgfdecorationsegmentamplitude+%
    \pgfdecorationsegmentaspect\pgfdecorationsegmentamplitude to last,
               width=+\pgfdecorationsegmentlength]
    { % line in the back
    %
    \pgfsetstrokecolor{\TDCoilColor}
    \pgfsetfillcolor{\TDCoilColor}
    \pgfsetstrokeopacity{\TDCoilOpacity}
    \pgfpathmoveto{\pgfpointanchor{TD@coilast}{center}}
    \pgfsetlinewidth{\TDCoilWidth} 
    \pgfpathcurveto
    {\pgfpoint@onthreedcoil{1.555}{ 1    }{4}}
    {\pgfpoint@onthreedcoil{2    }{ 0.555}{5}}
    {\pgfpoint@onthreedcoil{2    }{ 0    }{6}}
    \pgfpathcurveto
    {\pgfpoint@onthreedcoil{2    }{-0.555}{7}}
    {\pgfpoint@onthreedcoil{1.555}{-1    }{8}}
    {\pgfpoint@onthreedcoil{1    }{-1    }{9}}
    \pgfusepath{stroke} 
    %
    % white background for front thick part
    %
    % \pgfsetstrokeopacity{1}
    % \pgfsetstrokecolor{white}
    % \pgfsetfillcolor{white}
    % \pgfsetlinewidth{1.5*\TDCoilWidth+1.5*\TDCoilDist}
    % \pgfpathmoveto{\pgfpoint@onthreedcoil{1    }{ 1    }{3}}
    % \pgfpathmoveto{\pgfpoint@onthreedcoil{1    }{-1    }{9}}
    % % draw forward
    % \pgfpathcurveto
    % {\pgfpoint@onthreedcoil{0.445}{-1    }{10}}
    % {\pgfpoint@onthreedcoil{0    }{-0.555}{11.25}}
    % {\pgfpoint@onthreedcoil{0    }{ 0    }{12.5}}
    % \pgfpathcurveto
    % {\pgfpoint@onthreedcoil{0    }{ 0.555}{13.25}}
    % {\pgfpoint@onthreedcoil{0.445}{ 1    }{14.25}}
    % {\pgfpoint@onthreedcoil{1    }{ 1    }{15}}
    % % draw the curve back
    % \pgfpathcurveto
    % {\pgfpoint@onthreedcoil{0.445}{ 1    }{14}}
    % {\pgfpoint@onthreedcoil{0    }{ 0.555}{12.75}}
    % {\pgfpoint@onthreedcoil{0    }{ 0    }{11.5}}
    % \pgfpathcurveto
    % {\pgfpoint@onthreedcoil{0    }{-0.555}{10.75}}
    % {\pgfpoint@onthreedcoil{0.445}{-1    }{10}}
    % {\pgfpoint@onthreedcoil{1    }{-1    }{9}}
    % \pgfusepath{stroke,fill} 
    % % 
    % % draw the thick foreground path
    % %
    % \pgfsetstrokecolor{\TDCoilColor}
    % \pgfsetfillcolor{\TDCoilColor}
    % \pgfsetstrokeopacity{\TDCoilOpacity}
    % \pgfpathmoveto{\pgfpoint@onthreedcoil{1    }{ 1    }{3}}
    % \pgfsetlinewidth{\TDCoilWidth} 
    % % forward shifted +
    % \pgfpathmoveto{\pgfpoint@onthreedcoil{1    }{-1    }{9}}
    % \pgfpathcurveto
    % {\pgfpoint@onthreedcoil{0.445}{-1    }{10}}
    % {\pgfpoint@onthreedcoil{0    }{-0.555}{11.25}}
    % {\pgfpoint@onthreedcoil{0    }{ 0    }{12.5}}
    % \pgfpathcurveto
    % {\pgfpoint@onthreedcoil{0    }{ 0.555}{13.25}}
    % {\pgfpoint@onthreedcoil{0.445}{ 1    }{14.25}}
    % {\pgfpoint@onthreedcoil{1    }{ 1    }{15}}
    % % draw the curve back shfted -
    % \pgfpathcurveto
    % {\pgfpoint@onthreedcoil{0.445}{ 1    }{14}}
    % {\pgfpoint@onthreedcoil{0    }{ 0.555}{12.75}}
    % {\pgfpoint@onthreedcoil{0    }{ 0    }{11.5}}
    % \pgfpathcurveto
    % {\pgfpoint@onthreedcoil{0    }{-0.555}{10.75}}
    % {\pgfpoint@onthreedcoil{0.445}{-1    }{10}}
    % {\pgfpoint@onthreedcoil{1    }{-1    }{9}}
    % \pgfusepath{stroke,fill} % <- added
    \pgfcoordinate{TD@coilast}{\pgfpoint@onthreedcoil{1    }{ 1    }{15}} 
  }
  \state{last}[width=.25\pgfdecorationsegmentlength+%
    \pgfdecorationsegmentaspect\pgfdecorationsegmentamplitude+%
    \pgfdecorationsegmentaspect\pgfdecorationsegmentamplitude,next state=final]
  {
    \pgfsetstrokecolor{\TDCoilColor}
    \pgfsetstrokeopacity{\TDCoilOpacity}
    \pgfsetlinewidth{\TDCoilWidth}
    \pgfpathmoveto{\pgfpointanchor{TD@coilast}{center}}
    \pgfpathcurveto
    {\pgfpoint@onthreedcoil{1.555}{ 1    }{4}}
    {\pgfpoint@onthreedcoil{2    }{ 0.555}{5}}
    {\pgfpoint@onthreedcoil{2    }{ 0    }{6}}
  }
  \state{final}
  {
    % \pgfpathlineto{\pgfpointdecoratedpathlast}
    \pgfusepath{stroke}
    \tikzset{restore stroke color/.expand once=\pgfsavedstrokecolor}
  }
}
\makeatother


%%%%%%%%%%%%%%%%%%%%
%%%%%%%%%%%%%%%%%%%
%%%%%%%%%%%%%%%%%%
% next coil


\pgfkeys{%
/pgf/decoration/.cd,
3d coil1 color/.store in=\TDCoilColor, 
3d coil1 color/.initial=black,
3d coil1 color=black,
3d coil1 width/.store in=\TDCoilWidth, 
3d coil1 width/.initial=0.4pt,
3d coil1 width=0.4pt,
3d coil1 dist/.store in=\TDCoilDist, 
3d coil1 dist/.initial=0.6pt,
3d coil1 dist=0.6pt,
3d coil1 opacity/.store in=\TDCoilOpacity, 
3d coil1 opacity/.initial=1,
3d coil1 opacity=1
}

\makeatletter % https://tex.stackexchange.com/a/219088/121799
\tikzset{get stroke color/.code={%
    \expandafter\global% Jump over, now we have \global
    \expandafter\let% Jump over now we have \global\let
    \expandafter\pgfsavedstrokecolor% Jump we have \global\let\pgf...
    \csname\string\color@pgfstrokecolor\endcsname% Finally expand this and put it at the end 
    },                                           % \global\let\pgf...{} in expanded form 
    restore stroke color/.code={\pgf@setstrokecolor#1},
}
\def\pgfpoint@onthreedcoil#1#2#3{%
  \pgf@x=#1\pgfdecorationsegmentamplitude%
  \pgf@x=\pgfdecorationsegmentaspect\pgf@x%
  \pgf@y=#2\pgfdecorationsegmentamplitude%
  \pgf@xa=0.083333333333\pgfdecorationsegmentlength%
  \advance\pgf@x by#3\pgf@xa%
  \advance\pgf@x by-\initialoffset pt%
}

% coil decoration
%
% Parameters: \pgfdecorationsegmentamplitude, \pgfdecorationsegmentlength,

\pgfdeclaredecoration{3d coil1}{initial}
{ 
    \state{initial}[width=0.25*\pgfdecorationsegmentlength+\pgfdecorationsegmentaspect*\pgfdecorationsegmentamplitude,
    next state=coil, persistent precomputation={
    \tikzset{get stroke color}
    \pgfmathsetmacro{\initialoffset}{0.25*\pgfdecorationsegmentlength+\pgfdecorationsegmentaspect*\pgfdecorationsegmentamplitude}
  }]
  {%
    \pgfpathmoveto{\pgfpointorigin} 
    \pgfsetstrokecolor{\TDCoilColor}
    \pgfsetstrokeopacity{\TDCoilOpacity}
    \pgfsetlinewidth{1.5*\TDCoilWidth} 
    \pgfpathcurveto
    {\pgfpoint@oncoil{0    }{ 0.555}{1}}
    {\pgfpoint@oncoil{0.445}{ 1    }{2}}
    {\pgfpoint@oncoil{1    }{ 1    }{3}}
    \pgfusepath{stroke}
    \pgfcoordinate{TD@coilast}{\pgfpoint@oncoil{1    }{ 1    }{3}} 
    }
  \state{coil}[switch if less than=%
    1.25\pgfdecorationsegmentlength+%
    \pgfdecorationsegmentaspect\pgfdecorationsegmentamplitude+%
    \pgfdecorationsegmentaspect\pgfdecorationsegmentamplitude to last,
               width=+\pgfdecorationsegmentlength]
    { % line in the back
    %
    % \pgfsetstrokecolor{\TDCoilColor}
    % \pgfsetfillcolor{\TDCoilColor}
    % \pgfsetstrokeopacity{\TDCoilOpacity}
    % \pgfpathmoveto{\pgfpointanchor{TD@coilast}{center}}
    % \pgfsetlinewidth{\TDCoilWidth} 
    % \pgfpathcurveto
    % {\pgfpoint@onthreedcoil{1.555}{ 1    }{4}}
    % {\pgfpoint@onthreedcoil{2    }{ 0.555}{5}}
    % {\pgfpoint@onthreedcoil{2    }{ 0    }{6}}
    % \pgfpathcurveto
    % {\pgfpoint@onthreedcoil{2    }{-0.555}{7}}
    % {\pgfpoint@onthreedcoil{1.555}{-1    }{8}}
    % {\pgfpoint@onthreedcoil{1    }{-1    }{9}}
    % \pgfusepath{stroke} 
    % %
    % white background for front thick part
    %
    \pgfsetstrokeopacity{1}
    \pgfsetstrokecolor{white}
    \pgfsetfillcolor{white}
    \pgfsetlinewidth{1.5*\TDCoilWidth+1.5*\TDCoilDist}
    \pgfpathmoveto{\pgfpoint@onthreedcoil{1    }{ 1    }{3}}
    \pgfpathmoveto{\pgfpoint@onthreedcoil{1    }{-1    }{9}}
    % draw forward
    \pgfpathcurveto
    {\pgfpoint@onthreedcoil{0.445}{-1    }{10}}
    {\pgfpoint@onthreedcoil{0    }{-0.555}{11.25}}
    {\pgfpoint@onthreedcoil{0    }{ 0    }{12.5}}
    \pgfpathcurveto
    {\pgfpoint@onthreedcoil{0    }{ 0.555}{13.25}}
    {\pgfpoint@onthreedcoil{0.445}{ 1    }{14.25}}
    {\pgfpoint@onthreedcoil{.7    }{ 1    }{15}}% вроде из-за этого белый цвет находит на верхние части витков при рисовании жирных поверх
    % draw the curve back
    \pgfpathcurveto
    {\pgfpoint@onthreedcoil{0.445}{ 1    }{14}}
    {\pgfpoint@onthreedcoil{0    }{ 0.555}{12.75}}
    {\pgfpoint@onthreedcoil{0    }{ 0    }{11.5}}
    \pgfpathcurveto
    {\pgfpoint@onthreedcoil{0    }{-0.555}{10.75}}
    {\pgfpoint@onthreedcoil{0.445}{-1    }{10}}
    {\pgfpoint@onthreedcoil{1    }{-1    }{9}}
    \pgfusepath{stroke,fill} 
    % 
    % draw the thick foreground path
    %
    \pgfsetstrokecolor{\TDCoilColor}
    \pgfsetfillcolor{\TDCoilColor}
    \pgfsetstrokeopacity{\TDCoilOpacity}
    \pgfpathmoveto{\pgfpoint@onthreedcoil{1    }{ 1    }{3}}
    \pgfsetlinewidth{\TDCoilWidth} 
    % forward shifted +
    \pgfpathmoveto{\pgfpoint@onthreedcoil{1    }{-1    }{9}}
    \pgfpathcurveto
    {\pgfpoint@onthreedcoil{0.445}{-1    }{10}}
    {\pgfpoint@onthreedcoil{0    }{-0.555}{11.25}}
    {\pgfpoint@onthreedcoil{0    }{ 0    }{12.5}}
    \pgfpathcurveto
    {\pgfpoint@onthreedcoil{0    }{ 0.555}{13.25}}
    {\pgfpoint@onthreedcoil{0.445}{ 1    }{14.25}}
    {\pgfpoint@onthreedcoil{1    }{ 1    }{15}}
    % draw the curve back shfted -
    \pgfpathcurveto
    {\pgfpoint@onthreedcoil{0.445}{ 1    }{14}}
    {\pgfpoint@onthreedcoil{0    }{ 0.555}{12.75}}
    {\pgfpoint@onthreedcoil{0    }{ 0    }{11.5}}
    \pgfpathcurveto
    {\pgfpoint@onthreedcoil{0    }{-0.555}{10.75}}
    {\pgfpoint@onthreedcoil{0.445}{-1    }{10}}
    {\pgfpoint@onthreedcoil{1    }{-1    }{9}}
    \pgfusepath{stroke,fill} % <- added
    \pgfcoordinate{TD@coilast}{\pgfpoint@onthreedcoil{1    }{ 1    }{15}} 
  }
  \state{last}[width=.25\pgfdecorationsegmentlength+%
    \pgfdecorationsegmentaspect\pgfdecorationsegmentamplitude+%
    \pgfdecorationsegmentaspect\pgfdecorationsegmentamplitude,next state=final]
  {
    \pgfsetstrokecolor{\TDCoilColor}
    \pgfsetstrokeopacity{\TDCoilOpacity}
    \pgfsetlinewidth{\TDCoilWidth}
    \pgfpathmoveto{\pgfpointanchor{TD@coilast}{center}}
    \pgfpathcurveto
    {\pgfpoint@onthreedcoil{1.555}{ 1    }{4}}
    {\pgfpoint@onthreedcoil{2    }{ 0.555}{5}}
    {\pgfpoint@onthreedcoil{2    }{ 0    }{6}}
  }
  \state{final}
  {
    % \pgfpathlineto{\pgfpointdecoratedpathlast}
    \pgfusepath{stroke}
    \tikzset{restore stroke color/.expand once=\pgfsavedstrokecolor}
  }
}
\makeatother




    \begin{figure}[H]
    \centering

    \begin{tikzpicture}[font=\small,
    >={Stealth[bend,inset=0pt,round,length=3mm, angle'=20]},vec/.style={>={Stealth[inset=0pt,round,length=3.5mm, angle'=20]}}]


\def\lm{2} 


% coil back
\draw [decoration={3d coil color=brown,aspect=.2, segment length=2.62mm, amplitude=5mm,3d coil,post length=0},decorate] (0,0) --++ (4.2,0);
\draw[brown,line cap=round] (4,0) --++ (0,-.5) --++ (2,0);


\begin{scope}[gray,line cap=round]
\def\dg{1.1} % должно сходиться с \dg в %3 линии
% \clip[] ({-\dg*1cm-\pgflinewidth},{{-2.52*1cm-4*\pgflinewidth}}) rectangle++ ({2*\lm*1cm+\dg*1cm+\dg*1cm+2*\pgflinewidth},{2*2.52*1cm+8*\pgflinewidth})
% (0,{-\lm/4}) rectangle++ ({2*\lm},{\lm/2});

%lines
%1
\def\ds{.25}
\foreach \i in {0,1,2}{
\draw
[yshift=-{(\lm*\i*1cm/4)*1/3}]
({.1*\lm*1cm},{(\lm/4)*2/3}) --++ ({1.8*\lm*1cm},0) 
coordinate (r\i)
; 
}

\path[] (r0) 
arc [start angle=-90, end angle=0, x radius={\ds*1cm}, y radius=.2cm] 
% arc [start angle=180, end angle=360, x radius={1cm+\ds*1cm}, y radius=.4cm] 
coordinate (rr0)
;

\draw[->] (rr0)
arc [start angle=0, end angle=90, x radius={1.8cm+\ds*1cm}, y radius=.5cm]
coordinate (e1)
;

\draw[] (e1)
arc [start angle=90, end angle=180, x radius={1.8cm+\ds*1cm}, y radius=.5cm]
arc [start angle=180, end angle=270, x radius={\ds*1cm}, y radius=.2cm] 
;



%2
\def\df{.7}
\path[] (r1) 
arc [start angle=-90, end angle=0, x radius={\df*1cm}, y radius=.5cm] 
% arc [start angle=180, end angle=360, x radius={1cm+\ds*1cm}, y radius=.4cm] 
coordinate (rr1)
;

\draw[->] (rr1)
arc [start angle=0, end angle=90, x radius={1.8cm+\df*1cm}, y radius=1.2cm]
coordinate (e1)
;

\draw[] (e1)
arc [start angle=90, end angle=180, x radius={1.8cm+\df*1cm}, y radius=1.2cm]
arc [start angle=180, end angle=270, x radius={\df*1cm}, y radius=.5cm] 
;


% %3
% \def\dg{1.1}
% \draw[yshift=.8cm,->] ({2*\lm*1cm+\dg*1cm},{(\lm/4)*7/8}) 
% arc [start angle=0, end angle=90, x radius={1cm+\dg*1cm}, y radius=1.5cm] 
% coordinate (e1)
% ;

% \draw[yshift=.8cm] (e1) 
% arc [start angle=90, end angle=180, x radius={1cm+\dg*1cm}, y radius=1.5cm] 
% arc [start angle=180, end angle=360, x radius={1cm+\dg*1cm}, y radius=1.1cm] 
% ;


%%%%%
%%%%%


%reverslines
\begin{scope}[yscale=-1]
%1
\def\ds{.25}
\foreach \i in {0,1}{
\draw
[yshift=-{(\lm*\i*1cm/4)*1/3}]
({.1*\lm*1cm},{(\lm/4)*2/3}) --++ ({1.8*\lm*1cm},0) 
coordinate (r\i)
; 
}

\draw[] (r0) 
arc [start angle=-90, end angle=0, x radius={\ds*1cm}, y radius=.2cm] 
% arc [start angle=180, end angle=360, x radius={1cm+\ds*1cm}, y radius=.4cm] 
coordinate (rr0)
;

\draw[->] (rr0)
arc [start angle=0, end angle=90, x radius={1.8cm+\ds*1cm}, y radius=.5cm]
coordinate (e1)
;

\draw[] (e1)
arc [start angle=90, end angle=180, x radius={1.8cm+\ds*1cm}, y radius=.5cm]
arc [start angle=180, end angle=270, x radius={\ds*1cm}, y radius=.2cm] 
;



%2
\def\df{.7}
\draw[yshift=.19cm] (r1) 
arc [start angle=-90, end angle=0, x radius={\df*1cm}, y radius=.5cm] 
% arc [start angle=180, end angle=360, x radius={1cm+\ds*1cm}, y radius=.4cm] 
coordinate (rr1)
;

\draw[->] (rr1)
arc [start angle=0, end angle=90, x radius={1.8cm+\df*1cm}, y radius=1.2cm]
coordinate (e1)
;

\draw[] (e1)
arc [start angle=90, end angle=180, x radius={1.8cm+\df*1cm}, y radius=1.2cm]
arc [start angle=180, end angle=270, x radius={\df*1cm}, y radius=.5cm] 
;
\end{scope}
\end{scope}





%magnet
% \fill[red,semitransparent] (0,{-\lm/4}) rectangle++ (\lm,{\lm/2});
% \fill[NavyBlue,semitransparent] (\lm,{-\lm/4}) rectangle++ (\lm,{\lm/2});


% \draw (0,{-\lm/4}) rectangle++ ({2*\lm},{\lm/2});


%%%%%%%%%%%%%%%%%%
%%%%%%%%%%%%%%%%%%
%%%%%%%%%%%%%%%%%%
% coil front

% \draw [decoration={3d coil color=brown,aspect=.2, segment length=2.62mm, amplitude=5mm,3d coil,post length=0},decorate] (0,0) --++ (4.2,0);
\draw [decoration={3d coil color=brown,aspect=.2, segment length=2.62mm, amplitude=5mm,3d coil1,post length=0},decorate] (0,0) --++ (4.2,0);
% \draw[blue,line cap=round] (0,1.86) --++ (1,0) (0,0) --++ (1,0);

\begin{scope}[line cap=round]
    
% доп лин
\def\ds{.25}
\foreach \i in {0,1,2}{
\path
[yshift=-{(\lm*\i*1cm/4)*1/3}]
({.1*\lm*1cm},{(\lm/4)*2/3}) --++ ({1.8*\lm*1cm},0) 
coordinate (r\i)
; 
}

\draw[gray] (r0) 
arc [start angle=-90, end angle=0, x radius={\ds*1cm}, y radius=.2cm] 
% arc [start angle=180, end angle=360, x radius={1cm+\ds*1cm}, y radius=.4cm] 
coordinate (rr0)
;

\def\df{.7}
\draw[gray] (r1) 
arc [start angle=-90, end angle=0, x radius={\df*1cm}, y radius=.5cm] 
% arc [start angle=180, end angle=360, x radius={1cm+\ds*1cm}, y radius=.4cm] 
coordinate (rr1)
;


\draw[gray,decoration={markings,mark=at position {1/2+1.5/10} with {\arrow{<}}},postaction={decorate}] (0.2,0) --++ (-1,0); 
\draw[gray,decoration={markings,mark=at position {1/2+1.5/10} with {\arrow{>}}},postaction={decorate}] ({1.9*\lm},0) --++ (1,0);


%I
\draw[brown,semithick] (0,{\pgflinewidth}) --++ (0,{-.5cm-\pgflinewidth}) --++ (-2,0) coordinate (I);
\draw[red,<-] ([yshift=-.2cm]I) --++ (1,0) node[midway,black,below] {$I$};

\end{scope}



    \end{tikzpicture}
\caption{Линии магнитного поля катушки с током}
\label{fig:p151}
\end{figure}

Линии поля катушки как бы выходят из нее со одной стороны и входят в нее с другой стороны (внутри катушки линии замыкаются).

Для определения того, с какой стороны из катушки выходят линии ее магнитного поля, можно воспользоваться следующим правилом.
%
\begin{law}
    \textbf{Правило правой руки.} Направление линий магнитного поля \emph{внутри} катушки укажет большой палец правой руки, если остальные согнутые пальцы (охватывающие катушку) расположены так, что они указывают направление тока в катушке.
\end{law}
%
Рисунок~\ref{fig:p152} иллюстрирует применение этого правила.  

    \begin{figure}[H]
    \centering
\includegraphics[]{el/38coil_pic.pdf}%
\caption{Правило правой руки}
\label{fig:p152}
\end{figure}


Сопоставляя рисунки~\ref{fig:p151} и~\ref{fig:p152}, можно убедиться, что большой палец на рис.\,\ref{fig:p152} указывает направление линий магнитного поля внутри катушки.

При плотной намотке витков \emph{внутри} достаточно длинной катушки вдали от ее краев магнитное поле является \emph{однородным}: из рис.\,\ref{fig:p151} видно, что \emph{линии однородного поля являются параллельными прямыми одинаковой густоты.} 

























 
\end{document}